\section{Funzioni di Ordine Superiore e Array (Lab 16)}

\subsection{Funzioni come Valori}
In JS le funzioni sono "Cittadini di prima classe". Significa che puoi metterle dentro una variabile, passarle come argomenti o ritornarle.
Una **Funzione di Ordine Superiore** è proprio una funzione che accetta un'altra funzione come parametro.

\subsection{Metodi degli Array (Il "Coltellino Svizzero")}
Ecco i metodi fondamentali per lavorare con gli array senza usare il ciclo for.

\subsubsection{1. forEach (Esegui per tutti)}
Esegue del codice per ogni elemento. Non ritorna nulla.
\begin{lstlisting}[language=JavaScript]
let nomi = ["Mario", "Luigi"];
nomi.forEach((nome) => console.log("Ciao " + nome));
\end{lstlisting}

\subsubsection{2. map (Trasforma)}
\begin{concept}{Visualizzazione: Map}
Crea un **nuovo array** trasformando ogni elemento uno a uno.
\begin{verbatim}
   Input:  [ 1 , 2 , 3 ]
              |   |   |
             *2  *2  *2
              |   |   |
   Output: [ 2 , 4 , 6 ]
\end{verbatim}
\end{concept}

\begin{jsbox}
let num = [1, 2, 3];
let doppi = num.map(x => x * 2); // [2, 4, 6]
\end{jsbox}

\subsubsection{3. filter (Seleziona)}
\begin{concept}{Visualizzazione: Filter}
Crea un **nuovo array** solo con gli elementi che passano il test (return true).
\begin{verbatim}
   Input:  [ 10 , 5 , 20 ]
              |   |    |
             OK   NO  OK
              |        |
   Output: [ 10 ,     20 ]
\end{verbatim}
\end{concept}

\begin{jsbox}
let num = [10, 5, 20];
let grandi = num.filter(x => x > 9); // [10, 20]
\end{jsbox}

\subsubsection{4. find e findIndex (Cerca)}
\begin{concept}{Differenza Find vs FindIndex}
\texttt{find} ti dà il \textbf{valore} del primo elemento trovato. \texttt{findIndex} ti dà la sua \textbf{posizione}.
\end{concept}
\begin{jsbox}
let num = [5, 12, 8];
let trovato = num.find(x => x > 10); // 12
let indice = num.findIndex(x => x > 10); // 1
\end{jsbox}

\subsubsection{5. every e some (Verifica)}
\begin{concept}{Verifica Booleana}
Restituiscono un Booleano (True/False).
\begin{itemize}
    \item \textbf{every:} Tutti soddisfano la condizione?
    \item \textbf{some:} Almeno uno soddisfa la condizione?
\end{itemize}
\end{concept}
\begin{lstlisting}[language=JavaScript]
let num = [2, 4, 6];
let tuttiPari = num.every(x => x % 2 == 0); // true
\end{lstlisting}

\subsubsection{6. reduce (Accumula)}
Riduce tutto l'array a un singolo valore (es. somma).
\begin{lstlisting}[language=JavaScript]
let num = [1, 2, 3];
let somma = num.reduce((acc, x) => acc + x, 0); // 6
\end{lstlisting}

\subsubsection{7. sort (Ordina)}
Ordina l'array **modificandolo** (In-Place).
Attenzione: di base ordina come stringhe! Per i numeri serve una funzione di comparazione.
\begin{lstlisting}[language=JavaScript]
let numeri = [10, 2, 5];
numeri.sort((a, b) => a - b); // [2, 5, 10] (Crescente)
numeri.sort((a, b) => b - a); // [10, 5, 2] (Decrescente)
\end{lstlisting}

\subsection{Closures (I Ricordi delle Funzioni)}
Una Closure è quando una funzione figlia "ricorda" le variabili della funzione madre, anche dopo che la madre ha finito.

\begin{lstlisting}[language=JavaScript]
function creaContatore() {
    let count = 0; // Variabile nascosta nella closure
    return function() {
        count++; // La funzione interna puo' ancora vederla!
        return count;
    }
}
const c = creaContatore();
console.log(c()); // 1
console.log(c()); // 2
\end{lstlisting}
